\section{A penalty reserve does not repair coarse spectral support replacement}
\label{sec:op3-spectral-load-reserve-obstruction}

One might try to avoid the squared-accuracy spectral requirement by
decreasing the obstacle penalty, thereby enlarging the approximate
matrix's support. The next construction rules out this generic support
rule even with a reserve polynomial in the logarithm of inverse accuracy.
It does not concern accurate refinement, and trees already admit fast
local algorithms in this note.

\begin{theorem}[Finite binary-tree obstruction with a penalty reserve;
\statusproved]
\label{thm:op3-spectral-load-reserve-obstruction}
Put $q=3/17$, $K=17/31$ and $\beta=289/300$. For any integer
$R\geq64$, take the original complete binary tree of depth $D=4R$,
rooted and seeded at $v$. Set
\[
 \alpha=1/3,\quad t=\bar\alpha=1/2,\quad
 \lambda=Kq^R/4,\quad e=2\lambda,\quad\delta=e/8.
\]
The graph has maximum degree three and volume greater than $4/e$.
Let $M=D_G-A/2$, using $D_G$ for the original degree matrix, and
let $L_{\rm right}$ contain just the original edges to right children.
The matrix $\widehat M=M+L_{\rm right}/4$ satisfies
\[
 M\preceq\widehat M\preceq\tfrac32M,\qquad
 \widehat M\one=M\one=d/2,
\]
and is a strictly diagonally dominant M-matrix. At the all-left vertex
$w$ of depth $R$, the original obstacle at load $e_v-\lambda d$ has
$u_w\geq2\lambda>\delta$. But the $\widehat M$ obstacle at load
$e_v-(\lambda/C)d$ has $\widehat u_{C,w}=0$ whenever
\begin{equation}
 C\geq1,\qquad C\beta^R\leq51/403.
 \label{eq:op3-reserve-failure-condition}
\end{equation}
This includes every fixed reserve, and every reserve polynomial in
$\log(2+1/e)$ for all sufficiently large $R$. Therefore constant spectral
quality and a polylogarithmic penalty reserve do not by themselves give
a significant-envelope finder, even at fixed positive teleportation.
\end{theorem}
\begin{proof}
The volume is $2^{4R+2}-4>4/e$. For example, it is at least
$2^{4R+1}$, and its ratio to $4/e=(8/K)(17/3)^R$ is at least
$(K/4)(48/17)^R>1$ for $R\geq2$.
Consider the depth-dependent nonnegative vector
\[
 x_i=\bigl(Kq^{\operatorname{depth}(i)}-2\lambda\bigr)_+.
\]
It is positive precisely through depth $R$, and $x_w=2\lambda$.
At positive nonroot coordinates the untruncated homogeneous coefficient is
\[
 3-\frac{1/2}{q}-q=-1/102<0.
\]
The constant term contributes $-3\lambda$, so the original row has
$(Mx)_i\leq-3\lambda$. Clipping a negative child value to zero only
decreases this row of $Mx$. At the root, $K(2-q)=1$ gives
$(Mx)_v=1-2\lambda$. Thus $Mx\leq e_v-\lambda d$ on positive
coordinates of $x$. The obstacle maximum principle gives $x\leq u$.
All these coordinates lie above the true degree-one leaves, since $D=4R$.

Since $0\preceq L_{\rm right}\preceq L$ and
$M=(L+D_G)/2$, the stated spectral sandwich follows. The perturbation
has zero row sums. Its conductances are $1/2$ on left edges and $3/4$
on right edges, while its grounding at vertex $i$ remains $d_i/2$.

Let $g=\widehat M^{-1}e_v$. We bound $g$ by finite-tree elimination,
including its true leaves. A descendant subtree with effective ground
conductance $H$, excluding the incoming edge, transfers a parent
potential through an edge of conductance $c$ by the factor $c/(c+H)$.
This follows by solving its one retained scalar equation. At a true leaf
$H_0=1/2$. At an internal nonroot vertex whose child subtrees have the
same remaining depth, the effective ground is
\[
 F(H)=\frac32+\frac{(1/2)H}{1/2+H}+\frac{(3/4)H}{3/4+H}.
\]
The function is increasing, and direct rational evaluation gives
\[
 H_1=F(H_0)=41/20,\qquad
 H_2=F(H_1)=\frac32+\frac{41}{102}+\frac{123}{224}
 =\frac{28001}{11424}>\frac{83}{34}.
\]
Since $H_1>H_0$, all subtrees of remaining depth at least two have
$H\geq H_2$. Every left transfer along the root-to-$w$ path is therefore
less than $17/100$. The root grounding equals one, so $g_v\leq1$,
and consequently
\[
 g_w\leq(17/100)^R.
\]
No infinite graph or uncharged numerical primitive is used in this bound;
the Schur equations are comparison identities on the specified finite tree.

Because $g\geq0$ and $\widehat Mg=e_v\geq e_v-(\lambda/C)d$,
the maximum principle gives $\widehat u_C\leq g$.
The vertex $w$ has original degree three and diagonal
$\widehat M_{ww}=13/4$: its incoming and left-child conductances are
$1/2$, its right-child conductance is $3/4$, and its grounding is $3/2$.
If $\widehat u_{C,w}>0$, stationarity and the zero seed load at $w$
would require
\[
 \frac{13}{4}\widehat u_{C,w}
 =\sum_{j\sim w}c_{wj}\widehat u_{C,j}-\frac{3\lambda}{C}
 \leq\sum_{j\sim w}c_{wj}g_j-\frac{3\lambda}{C}
 =\frac{13}{4}g_w-\frac{3\lambda}{C}.
\]
But \cref{eq:op3-reserve-failure-condition} is exactly the sufficient
inequality $(13/4)(17/100)^R\leq3Kq^R/(4C)$ making the right side
nonpositive. Positivity is impossible, proving omission of $w$.

For an explicit uniform sequence, take $C=2^k$ and $R=64(k+1)$.
The integer inequality $8\cdot289^{64}<300^{64}$ gives
$C\beta^R<2^k8^{-(k+1)}\leq1/8<51/403$.
For any fixed reserve choose a larger dyadic one; increasing the penalty
can only decrease obstacle coordinates. Finally
$\log(1/e)=R\log(17/3)+\log(2/K)$, while $\beta^R$ decays
exponentially. Every fixed polynomial in that logarithm satisfies
\cref{eq:op3-reserve-failure-condition} eventually.
\end{proof}

\paragraph{Audit and interpretation.}
\path{spectral_load_reserve_obstruction.py} checks the rational subsolution,
the two-step ground invariant, and 257 exact dyadic reserve/volume
certificates through $C=2^{256}$. These last trees are symbolic objects;
they are not materialized or claimed as executed local solves. Small
independent dense computations check five finite-tree Green functions
on 243 total vertices, every Schur/lifting identity, ten left-transfer
bounds, and five original radial obstacle quotients. The quotient checks
include 129 exact KKT and subsolution rows. Seven further witnesses use
reserves polynomial in $R+1$.
Hashes, exact parameters and measured audit duration are recorded in
\path{SPECTRAL_LOAD_RESERVE_OBSTRUCTION_AUDIT.json}.

This refutes the implication from the stated coarse spectral guarantee
and a penalty reserve to support containment. It does not assert that any
particular sparsifier returns the constructed matrix, nor does it constrain
a method that verifies/refines original residuals or deliberately explores
additional coordinates. It is not an OP3 computational lower bound.
